% 夜半用户手册 · 第 6 章：自动化与调制（中文）
% T-DOC-001 逐章交付（CH06 任务未到前预写）。
% 事实来源：yeban-model/src/project.rs 的 AutomationWriteMode、bridge.rs 自动化泳道、
% cli.rs SHORTCUTS、yeban-dsp/src/oscillator.rs 的全局 LFO、docs/manual/05-mixer-effects.md。

\chapter{自动化与调制}
\label{ch:automation-modulation}

\begin{quote}
易穷则变，变则通，通则久。
\end{quote}

出处：《周易·系辞下》

自动化让参数随时间变化；懂得变通，声音才不会僵住。

自动化让参数随时间变化：你可以画一条曲线，也可以边播放边把操作录进去。
这一章讲怎么显示泳道、四种写入模式各是什么意思，以及当前版本在调制上的
真实边界。

\section{自动化曲线与写入模式}

\subsection{显示泳道}

在编曲区按 A，或用「视图」里的自动化开关（Cmd/Ctrl+Shift+A）显示自动化泳道。
每条非主总线的轨可以有自己的泳道，按目标参数排列；泳道里是一条折线，顶点
按时间顺序连接，横轴与编曲区的时间对齐，因此可以直接对照片段位置来读曲线。

\screenshot{ch06-automation-lanes}{编曲区展开的自动化泳道，显示一条带顶点的参数曲线}

\subsection{画曲线与改顶点}

泳道里的折线由若干顶点组成。添加顶点即在某个时间点给参数定一个值，两个
顶点之间按折线插值。拖动顶点可以改变它的时间或数值；删除顶点则让曲线在
相邻两点之间重新连起来。曲线随工程保存，播放时按这条曲线驱动参数，导出时
同样生效。

\subsection{四种写入模式}

夜半的模型定义了四种自动化写入模式，理解它们能避免「录了一堆没用的曲线」：

\begin{description}
  \item[Off（不录制）] 默认值。播放时不写入任何自动化；手动改参数不会改动曲线。
  \item[Write（写到底）] 走带时只要参数被改动就持续写入，直到你显式停止。
  \item[Touch（触碰写入）] 只在你正手动触碰该参数时写入，松手即停。适合「推一下再放开」的修正。
  \item[Latch（锁存写入）] 一旦开始触碰就持续写入，松手也不停，直到显式停止。适合把一段推子动作整段锁进工程。
\end{description}

\begin{table}[htbp]
\centering
\caption{四种写入模式对照}
\label{tab:ch06-write-modes}
\begin{tabular}{ll}
\hline
模式 & 松手后的行为 \\
\hline
Off & 不写入 \\
Write & 持续写（需显式停止） \\
Touch & 立即停 \\
Latch & 继续写（需显式停止） \\
\hline
\end{tabular}
\end{table}

\section{参数绑定与 LFO}

\subsection{参数绑定}

自动化泳道绑定的是具体轨道上的具体参数：选择目标参数后就在这条泳道上画。
设备参数（如插入效果的参数）也可以成为自动化目标；把设备链插好后，在泳道
的目标里选对应参数即可。参数改动与自动化写入共用同一条提交路径，因此撤销
可以逐步还原，不会出现「曲线改了但工程没记」的情况。

\subsection{LFO 的真实边界}

夜半的 DSP 核心里确实实现了一个\textbf{全局 LFO}（低频振荡器）：所有声部共用
一份，按块求值，波形可选。但要特别注意：目前它\textbf{还没有接进合成器的
调制入口}——也就是说，不要指望在界面上打开一个 LFO 旋钮就得到颤音或滤波
摆动；这些调制路径仍在演进中。

\begin{figure}[htbp]
\centering
\begin{tikzpicture}[node distance=1.1cm, >=stealth,
    box/.style={draw, rounded corners, align=center, minimum width=2.3cm, minimum height=0.8cm}]
  \node[box] (a) {选自动化目标};
  \node[box, right of=a] (b) {选写入模式};
  \node[box, right of=b] (c) {画曲线/边播边写};
  \node[box, right of=c] (d) {播放按曲线驱动};
  \node[box, below of=c] (e) {LFO：DSP 已有，未接合成器};
  \draw[->] (a)--(b); \draw[->] (b)--(c); \draw[->] (c)--(d);
  \draw[dashed,->] (c)--(e);
\end{tikzpicture}
\caption{自动化流程与 LFO 边界：实线为当前可用路径，虚线为尚未接入调制的路径}
\label{fig:ch06-automation-flow}
\end{figure}

因此，本章不把「颤音」「滤波摆动」写成已提供的功能；如果你在界面上找不到
对应控件，那就是它还没接上，不是你漏看了。

\section{帮助主题映射}

\begin{table}[htbp]
\centering
\caption{本章小节 $\to$ 应用内帮助锚点}
\label{tab:ch06-help}
\begin{tabular}{ll}
\hline
本章小节 & 帮助主题 \\
\hline
显示泳道 & \texttt{automation-show-lanes} \\
画曲线与改顶点 & \texttt{automation-edit-curve} \\
四种写入模式 & \texttt{automation-write-modes} \\
参数绑定 & \texttt{automation-target-binding} \\
LFO 的真实边界 & \texttt{modulation-lfo-boundary} \\
\hline
\end{tabular}
\end{table}

\paragraph{本章学完能做什么} 显示并编辑自动化曲线，说清 Off/Write/Touch/Latch 的区别，并知道 LFO 目前只有 DSP 核心实现、尚未接入合成器。

\paragraph{动手练习} 在一条轨上展开自动化泳道，选音量参数，用 Touch 模式边播放边推一下音量再松手，确认只写入触碰那一段；改回 Off，手动拖动一个顶点并试听差异；最后找一找界面上有没有 LFO 颤音控件，确认它当前尚未提供。
